% =============================================================================
% main.tex
% Guia do Usuário — latex-abnt
%
% Este documento funciona simultaneamente como:
%   1. Manual de referência de todos os comandos do projeto
%   2. Exemplo compilável de um trabalho acadêmico ABNT válido
%   3. Checklist de conformidade ABNT
%
% Compilação (a partir do diretório deste arquivo):
%   TEXINPUTS=../../src: latexmk -pdf -shell-escape main.tex
%
% Ou a partir da raiz do projeto (requer -output-directory):
%   TEXINPUTS=./src: latexmk -pdf -shell-escape -output-directory=templates/guia-usuario templates/guia-usuario/main.tex
% =============================================================================

\documentclass{abnt-trabalho-academico}

% Pacotes não carregados automaticamente pela classe:
\usepackage{abnt-resumo}               % resumo e abstract
\usepackage{abnt-citacoes}             % citações (carrega abnt-refs automaticamente)
\usepackage{ibge-tabelas}              % tabelas no padrão IBGE
\usepackage{abnt-indice}               % índice remissivo
\usepackage{abnt-lombada}              % lombada

% O biblatex é carregado por abnt-refs (via abnt-citacoes).
\addbibresource{refs.bib}

% Pacote auxiliar para exemplos verbatim em minipage
\usepackage{fancyvrb}

% Comando \code para uso dentro de tabelaibge (onde \verb é proibido)
\newcommand{\code}[1]{\texttt{#1}}
\newcommand{\bsl}{\textbackslash}

% --- Metadados ---
\titulo{Guia do Usuário --- latex-abnt}
\subtitulo{classes e pacotes LaTeX para normas ABNT}
\autor{Equipe latex-abnt}
\orientador{(não aplicável)}
\instituicao{Projeto latex-abnt}
\local{Brasília}
\ano{2026}
\natureza{Guia de referência elaborado como exemplo de documento acadêmico
  formatado com as classes e pacotes do projeto latex-abnt.}
\programa{Documentação do projeto}

% Dados da lombada
\lombadaElementos{v.~1}
\lombadaLogomarcaTexto{ABNT}

% =============================================================================
\begin{document}

% =====================================================================
% PRÉ-TEXTUAIS
% =====================================================================
\imprimircapa
\imprimirfolhaderosto

% Folha de aprovação com placeholders (documento não é acadêmico real)
\dataaprovacao{1 de janeiro de 2026}
\membrodabanca{Colaborador A}{Desenvolvedor --- Projeto latex-abnt}
\membrodabanca{Colaborador B}{Revisor --- Projeto latex-abnt}
\membrodabanca{Colaborador C}{Testador --- Projeto latex-abnt}
\imprimirfolhadeaprovacao

\imprimirdedicatoria{%
  A todos os estudantes e pesquisadores brasileiros que dedicam
  tempo e esforço para formatar seus trabalhos conforme as normas ABNT.}

\imprimiragradecimentos{%
  Agradecemos à comunidade \LaTeX{} brasileira, em particular ao projeto
  abnTeX2, que inspirou este trabalho.

  Agradecemos também aos mantenedores do biblatex-abnt, do tabularray
  e de todos os pacotes dos quais este projeto depende.}

\imprimirepigrafe{%
  ``A tipografia é o ofício que dá forma visível e durável
  --- e portanto existência independente --- ao pensamento humano.''
  (Robert Bringhurst)}

\imprimirresumo{%
  Este documento apresenta o projeto latex-abnt, um conjunto de classes
  e pacotes LaTeX para formatação de documentos conforme normas da
  Associação Brasileira de Normas Técnicas. O guia descreve a instalação,
  configuração e uso de cada componente, cobrindo trabalhos acadêmicos
  e projetos de pesquisa. O próprio documento serve como exemplo compilável
  de um trabalho acadêmico ABNT, utilizando todos os recursos descritos.
  São apresentados os comandos para estrutura do documento, numeração
  progressiva de seções, citações e referências bibliográficas, tabelas
  no padrão IBGE, índice remissivo e lombada. O documento inclui ainda
  um checklist de conformidade com itens que as normas exigem mas que
  o LaTeX não automatiza, e uma referência rápida de todos os comandos
  públicos do projeto.}{%
  latex-abnt; ABNT; LaTeX; normas; trabalho acadêmico}

\imprimirabstract{%
  This document presents the latex-abnt project, a set of LaTeX classes
  and packages for formatting documents according to the standards of the
  Brazilian Association of Technical Standards (ABNT). The guide describes
  the installation, configuration and use of each component, covering
  academic works and research projects. The document itself serves as a
  compilable example of an ABNT academic work, using all the features
  described. Commands for document structure, progressive section numbering,
  citations and bibliographic references, IBGE-standard tables, back-of-book
  index and spine are presented. The document also includes a conformance
  checklist with items that standards require but LaTeX does not automate,
  and a quick reference of all public commands.}{%
  latex-abnt; ABNT; LaTeX; standards; academic work}

\imprimirsumario

% =====================================================================
% ELEMENTOS TEXTUAIS
% =====================================================================
\textual

% =====================================================================
% CAPÍTULO 1 — INTRODUÇÃO
% =====================================================================
\chapter{Introdução}\label{ch:introducao}

\index{introdução}

O projeto latex-abnt fornece classes e pacotes \LaTeX{} para formatar
documentos conforme as normas da Associação Brasileira de Normas Técnicas
(ABNT)\index{ABNT}. O objetivo é oferecer uma alternativa moderna e
modular ao abnTeX2, aproveitando os recursos do \LaTeX{} atual (formato
2020+, expl3).

Este guia cumpre três funções: manual de referência dos comandos,
exemplo compilável de trabalho acadêmico e checklist de conformidade.
O próprio documento que o leitor tem em mãos foi formatado com
\verb|\documentclass{abnt-trabalho-academico}| e demonstra na prática
todos os recursos descritos.

\section{Início rápido}\label{sec:inicio-rapido}

\index{início rápido}
Para produzir um documento ABNT funcional, são necessários três passos:

\begin{alineas}
  \item criar um arquivo \verb|main.tex| com o preâmbulo mínimo mostrado
    a seguir;
  \item preencher os metadados (\verb|\titulo|, \verb|\autor|, etc.)
    e o conteúdo;
  \item compilar com o comando indicado abaixo.
\end{alineas}

Preâmbulo mínimo:

\begin{Verbatim}[fontsize=\small]
\documentclass{abnt-trabalho-academico}
\usepackage{abnt-citacoes}
\addbibresource{refs.bib}

\titulo{Título do trabalho}
\autor{Nome do Autor}
\orientador{Nome do Orientador}
\instituicao{Nome da Instituição}
\local{Cidade}
\ano{2026}
\natureza{Dissertação apresentada ao Programa...}

\begin{document}
\imprimircapa
\imprimirfolhaderosto
\imprimirsumario
\textual

\chapter{Introdução}
Texto do trabalho \parencite{referencia}.

\postextual
\imprimirreferencias
\end{document}
\end{Verbatim}

Comando de compilação (a partir da raiz do projeto):

\begin{Verbatim}[fontsize=\small]
TEXINPUTS=./src: latexmk -pdf -shell-escape main.tex
\end{Verbatim}

O \verb|latexmk| detecta e invoca automaticamente o Biber (para
referências) e o makeindex (para índice remissivo). A opção
\verb|-shell-escape| é necessária porque o pacote \verb|abnt-indice|
usa \verb|imakeidx|, que chama \verb|makeindex| via \verb|\write18|.

\section{Requisitos}\label{sec:requisitos}

\index{requisitos}
O projeto requer:

\begin{alineas}
  \item TeX Live (pdfLaTeX\index{pdfLaTeX} ou
    XeLaTeX\index{XeLaTeX}) --- distribuição 2020 ou posterior;
  \item Biber\index{Biber} --- backend do biblatex;
  \item biblatex-abnt\index{biblatex-abnt} --- estilo ABNT para biblatex
    (disponível no CTAN e no TeX Live).
\end{alineas}

\section{Classes disponíveis}\label{sec:classes}

\index{classes}
O projeto oferece duas classes, ambas herdando de \verb|abnt.cls|:

\begin{alineas}
  \item \verb|abnt-trabalho-academico| --- para trabalhos acadêmicos
    conforme NBR 14724:2024 (teses, dissertações, TCCs);
  \item \verb|abnt-projeto-pesquisa| --- para projetos de pesquisa
    conforme NBR 15287:2025 (sem dedicatória, agradecimentos, epígrafe,
    folha de aprovação nem resumos).
\end{alineas}

Ambas as classes carregam automaticamente apenas \verb|abnt-sumario| e
\verb|abnt-numeracao|. Os demais pacotes devem ser incluídos pelo
usuário com \verb|\usepackage|.

\section{Pacotes e normas}\label{sec:pacotes-normas}

\index{pacotes}
A tabela a seguir resume qual pacote implementa qual norma:

\begin{tabela}{Pacotes do projeto latex-abnt e normas correspondentes}
  \label{tab:pacotes}
  \ibgeFonte{Elaboração própria.}
  \begin{tabelaibge}{colspec={llp{6cm}}}
    Pacote & Norma & Função \\
    \code{abnt-numeracao} & NBR 6024:2012 & Numeração progressiva de seções \\
    \code{abnt-sumario}   & NBR 6027:2012 & Formatação do sumário \\
    \code{abnt-resumo}    & NBR 6028:2021 & Resumo e abstract \\
    \code{abnt-indice}    & NBR 6034:2004 & Índice remissivo \\
    \code{abnt-refs}      & NBR 6023:2025 & Referências bibliográficas \\
    \code{abnt-citacoes}  & NBR 10520:2023 & Citações \\
    \code{ibge-tabelas}   & IBGE 1993     & Tabelas no padrão IBGE \\
    \code{abnt-lombada}   & NBR 12225:2023 & Lombada \\
  \end{tabelaibge}
\end{tabela}

% =====================================================================
% CAPÍTULO 2 — ESTRUTURA DO DOCUMENTO
% =====================================================================
\chapter{Estrutura do documento}\label{ch:estrutura}

\index{estrutura do documento}
Um trabalho acadêmico ABNT divide-se em três partes: pré-textuais,
textuais e pós-textuais. A classe \verb|abnt-trabalho-academico|
fornece comandos para cada elemento. A sequência recomendada é
apresentada a seguir.

\section{Metadados}\label{sec:metadados}

\index{metadados}
Os metadados são definidos no preâmbulo, antes de \verb|\begin{document}|.
Todos os comandos aceitam um argumento obrigatório:

\begin{Verbatim}[fontsize=\small]
\titulo{Título do trabalho}
\subtitulo{subtítulo, se houver}
\autor{Nome completo do autor}
\orientador{Nome do orientador}
\coorientador{Nome do coorientador}
\instituicao{Nome da instituição}
\local{Cidade}
\ano{2026}
\natureza{Dissertação apresentada ao Programa de...}
\programa{Programa de Pós-Graduação em...}
\area{Área de concentração}
\end{Verbatim}

Os comandos \verb|\titulo| e \verb|\autor| são obrigatórios para a capa
e folha de rosto. Os demais são opcionais (campos vazios são omitidos
silenciosamente). O hyperref preenche automaticamente os metadados PDF
(\verb|pdftitle|, \verb|pdfauthor|) a partir de \verb|\titulo| e
\verb|\autor|.

\section{Elementos pré-textuais}\label{sec:pretextuais}

\index{pré-textuais}
A sequência dos elementos pré-textuais segue a NBR 14724. O comando
\verb|\pretextual| é chamado automaticamente por \verb|\begin{document}|
(não é necessário invocá-lo manualmente).

\begin{Verbatim}[fontsize=\small]
\begin{document}
\imprimircapa
\imprimirfolhaderosto
% Folha de aprovação (após defesa):
\dataaprovacao{15 de março de 2026}
\membrodabanca{Prof. Dr. Fulano}{Doutor --- Universidade X}
\membrodabanca{Profa. Dra. Beltrana}{Doutora --- Universidade Y}
\imprimirfolhadeaprovacao
\imprimirdedicatoria{Texto da dedicatória.}
\imprimiragradecimentos{Texto dos agradecimentos.}
\imprimirepigrafe{``Citação inspiradora.'' (Autor)}
\imprimirresumo{Texto do resumo...}{palavra1; palavra2; palavra3}
\imprimirabstract{Abstract text...}{keyword1; keyword2; keyword3}
\imprimirsumario
\end{Verbatim}

Cada comando gera o elemento em página própria. A capa e a folha de
rosto são geradas por \verb|\imprimircapa| e \verb|\imprimirfolhaderosto|,
respectivamente.

A folha de aprovação requer \verb|\dataaprovacao{data}| e uma ou mais
chamadas de \verb|\membrodabanca{nome}{titulação}| antes de
\verb|\imprimirfolhadeaprovacao|. Se nenhum membro for definido,
placeholders são inseridos automaticamente.

O comando \verb|\imprimirdedicatoria{texto}| posiciona o texto na parte
inferior direita da página. O comando \verb|\imprimiragradecimentos{texto}|
aceita parágrafos múltiplos (argumento \verb|+m| no LaTeX3). O comando
\verb|\imprimirepigrafe{texto}| usa o mesmo layout da dedicatória.

\subsection{Paginação pré-textual}\label{subsec:paginacao-pre}

\index{paginação}
A paginação pré-textual segue a NBR 14724 seção 5.3: as páginas são
contadas em arábico desde a primeira página (capa), mas os números não
são exibidos. Quando \verb|\textual| é chamado, a contagem \emph{não}
reinicia --- a primeira página textual terá número maior que~1.

\section{Transição para o texto}\label{sec:textual}

O comando \verb|\textual| ativa a exibição dos números de página e
restaura o estilo \verb|fancy|. Deve ser chamado imediatamente antes
do primeiro capítulo:

\begin{Verbatim}[fontsize=\small]
\textual
\chapter{Introdução}
...
\end{Verbatim}

\section{Elementos pós-textuais}\label{sec:postextuais}

\index{pós-textuais}
O comando \verb|\postextual| marca o início dos elementos pós-textuais.
A paginação contínua é mantida. Os comandos típicos são:

\begin{Verbatim}[fontsize=\small]
\postextual
\imprimirreferencias   % gera a lista de referências
\imprimirindice        % gera o índice remissivo
\end{Verbatim}

O comando \verb|\imprimirreferencias| é um wrapper para
\verb|\printbibliography| com heading customizado (``REFERÊNCIAS'',
centralizado, em maiúsculas e negrito). O comando \verb|\imprimirindice|
aceita um argumento opcional \verb|[nome]| para índices nomeados
(quando se usa múltiplos índices via imakeidx).

\section{Projeto de pesquisa}\label{sec:projeto-pesquisa}

\index{projeto de pesquisa}
A classe \verb|abnt-projeto-pesquisa| implementa a NBR 15287:2025.
As diferenças em relação ao trabalho acadêmico são:

\begin{alineas}
  \item a capa é elemento opcional (não obrigatório);
  \item não existem dedicatória, agradecimentos, epígrafe, folha de
    aprovação nem resumos;
  \item o comando \verb|\postextual| \textbf{não existe} --- use
    \verb|\pagestyle{fancy}| diretamente se necessário.
\end{alineas}

Comandos disponíveis: \verb|\imprimircapa|, \verb|\imprimirfolhaderosto|,
\verb|\imprimirsumario|, \verb|\textual| e todos os metadados.
Comandos \textbf{não} disponíveis: \verb|\imprimirfolhadeaprovacao|,
\verb|\imprimirdedicatoria|, \verb|\imprimiragradecimentos|,
\verb|\imprimirepigrafe|, \verb|\postextual|.

Exemplo mínimo de projeto de pesquisa:

\begin{Verbatim}[fontsize=\small]
\documentclass{abnt-projeto-pesquisa}
\usepackage{abnt-citacoes}
\addbibresource{refs.bib}

\titulo{Título do projeto}
\autor{Nome do Autor}
\orientador{Nome do Orientador}
\instituicao{Universidade}
\local{Cidade}
\ano{2026}
\natureza{Projeto de pesquisa apresentado ao...}

\begin{document}
\imprimircapa
\imprimirfolhaderosto
\imprimirsumario
\textual

\chapter{Tema}
Texto do projeto \parencite{referencia}.

\imprimirreferencias
\end{document}
\end{Verbatim}


% =====================================================================
% CAPÍTULO 3 — NUMERAÇÃO E SEÇÕES
% =====================================================================
\chapter{Numeração e seções}\label{ch:numeracao}

\index{numeração progressiva}
\index{abnt-numeracao}
O pacote \verb|abnt-numeracao| implementa a NBR 6024:2012, que
regulamenta a numeração progressiva das seções de um documento. Este
pacote é carregado automaticamente pelas duas classes do projeto.

\section{Níveis de seção}\label{sec:niveis}

Com as classes ABNT (que herdam de \verb|report|), os cinco níveis
de seção são mapeados da seguinte forma:

\begin{alineas}
  \item nível primário: \verb|\chapter| --- MAIÚSCULAS + NEGRITO;
  \item nível secundário: \verb|\section| --- MAIÚSCULAS;
  \item nível terciário: \verb|\subsection| --- negrito;
  \item nível quaternário: \verb|\subsubsection| --- normal;
  \item nível quinário: \verb|\paragraph| --- itálico.
\end{alineas}

Este próprio guia demonstra a formatação: os capítulos estão em
maiúsculas e negrito, as seções em maiúsculas sem negrito, e assim
por diante. A hierarquia tipográfica segue a NBR 6024 seção 4.1.j.

\subsection{Exemplo de nível terciário}

Este é um exemplo de seção terciária (\verb|\subsection|), formatada
em negrito conforme a norma.

\subsubsection{Exemplo de nível quaternário}

Este é um exemplo de seção quaternária (\verb|\subsubsection|),
formatada em texto normal (sem negrito nem itálico).

\paragraph{Exemplo de nível quinário}

Este é um exemplo de seção quinária (\verb|\paragraph|), formatada
em itálico conforme a norma.

\section{Uso standalone com article}\label{sec:standalone-article}

\index{standalone}
Quando o pacote é usado com \verb|\documentclass{article}| (sem as
classes ABNT), a correspondência muda porque \verb|article| não possui
\verb|\chapter|:

\begin{alineas}
  \item nível primário: \verb|\section|;
  \item nível secundário: \verb|\subsection|;
  \item nível terciário: \verb|\subsubsection|;
  \item nível quaternário: \verb|\paragraph|;
  \item nível quinário: \verb|\subparagraph|.
\end{alineas}

\section{Títulos sem indicativo numérico}\label{sec:titulos-sem-numero}

\index{títulos sem indicativo}
Títulos sem indicativo numérico (como RESUMO, ABSTRACT, SUMÁRIO,
REFERÊNCIAS) são gerados com \verb|\chapter*{TÍTULO}|. O pacote
\verb|abnt-numeracao| formata esses títulos centralizados, com o mesmo
destaque tipográfico das seções primárias (maiúsculas e negrito),
conforme a NBR 6024 seção 4.1.h.

\section{Alíneas e subalíneas}\label{sec:alineas}

\index{alíneas}
\index{subalíneas}
O pacote fornece os ambientes \verb|alineas| e \verb|subalineas|,
conforme as seções 4.2 e 4.3 da NBR 6024:

\begin{Verbatim}[fontsize=\small]
\begin{alineas}
  \item primeiro item;
  \item segundo item:
    \begin{subalineas}
      \item primeiro subitem;
      \item segundo subitem;
    \end{subalineas}
  \item terceiro item.
\end{alineas}
\end{Verbatim}

Resultado:

\begin{alineas}
  \item primeiro item;
  \item segundo item:
    \begin{subalineas}
      \item primeiro subitem;
      \item segundo subitem;
    \end{subalineas}
  \item terceiro item.
\end{alineas}

As alíneas usam letras minúsculas seguidas de parêntese: a), b),~\ldots,
z), aa), bb),~\ldots, zz). O limite é de 52~itens; acima disso, um erro
é emitido. As subalíneas usam travessão como marcador.

A pontuação das alíneas (ponto e vírgula entre itens, ponto final no
último) é responsabilidade do usuário --- o \LaTeX{} não automatiza isso.

\section{Opção de maiúsculas}\label{sec:opcao-maiusculas}

\index{maiúsculas}
Por padrão, o pacote aplica maiúsculas automáticas nos títulos de seção
primária e secundária. Para desligar esse comportamento:

\begin{Verbatim}[fontsize=\small]
\usepackage[maiusculas=false]{abnt-numeracao}
\end{Verbatim}


% =====================================================================
% CAPÍTULO 4 — CITAÇÕES E REFERÊNCIAS
% =====================================================================
\chapter{Citações e referências}\label{ch:citacoes}

\index{citações}
\index{abnt-citacoes}
\index{abnt-refs}
Os pacotes \verb|abnt-citacoes| (NBR 10520:2023) e \verb|abnt-refs|
(NBR 6023:2025) configuram o biblatex-abnt para produzir citações e
referências conforme as normas ABNT. O \verb|abnt-citacoes| carrega
\verb|abnt-refs| automaticamente.

\section{Sistema autor-data}\label{sec:autor-data}

\index{autor-data}
O sistema autor-data é o padrão. Os principais comandos de citação são
fornecidos pelo biblatex-abnt:

\begin{alineas}
  \item \verb|\textcite{chave}| --- citação no corpo do texto:
    \textcite{abnt2024};
  \item \verb|\parencite{chave}| --- citação entre parênteses:
    \parencite{medeiros2019};
  \item \verb|\parencite[p.~42]{chave}| --- com indicação de página:
    \parencite[p.~42]{abnt2024}.
\end{alineas}

Com quatro ou mais autores, o biblatex-abnt aplica \textit{et al.}
automaticamente: \parencite{cervo2007}.

Quando há duas obras do mesmo autor no mesmo ano, o biblatex-abnt
desambigua com letras: \parencite{silva2020a} e \parencite{silva2020b}.

\section{Sistema numérico}\label{sec:numerico}

\index{sistema numérico}
Para usar o sistema numérico, altere a opção no preâmbulo:

\begin{Verbatim}[fontsize=\small]
\usepackage[sistema=numerico]{abnt-citacoes}
\end{Verbatim}

\section{Citação direta longa}\label{sec:citacao-direta}

\index{citação direta}
\index{citacaodireta}
Citações diretas com mais de três linhas devem usar o ambiente
\verb|citacaodireta|, que aplica recuo de 4~cm, fonte menor (10pt)
e espaçamento simples:

\begin{Verbatim}[fontsize=\small]
\begin{citacaodireta}
Texto da citação longa \parencite[p.~50]{chave}.
\end{citacaodireta}
\end{Verbatim}

Exemplo real:

\begin{citacaodireta}
A normalização de documentos acadêmicos no Brasil evoluiu significativamente
nas últimas décadas, acompanhando as transformações tecnológicas e as
demandas crescentes por padronização na produção científica. As normas
da ABNT estabelecem critérios claros para apresentação, citação e
referenciação, contribuindo para a qualidade e a acessibilidade do
conhecimento produzido nas universidades brasileiras
\parencite[p.~50]{medeiros2019}.
\end{citacaodireta}

A detecção automática de citações longas (mais de três linhas) não é
viável em \TeX{}; o usuário decide quando usar o ambiente.

\section{Citação de citação}\label{sec:apud}

\index{apud}
A citação de citação (apud) é feita combinando \verb|\textcite| ou
\verb|\parencite| com a indicação explícita da fonte intermediária.
O biblatex-abnt oferece suporte nativo a essa construção. Exemplo:

\begin{Verbatim}[fontsize=\small]
\textcite{autor-original} apud \textcite{autor-consultado}
\end{Verbatim}

\section{Referências}\label{sec:referencias}

\index{referências}
A lista de referências é gerada por \verb|\imprimirreferencias|,
que chama \verb|\printbibliography| com heading customizado. O título
``REFERÊNCIAS'' aparece centralizado, em maiúsculas e negrito, e é
adicionado ao sumário automaticamente.

\section{Backref}\label{sec:backref}

\index{backref}
O backref (indicação das páginas onde cada referência é citada) é
ativado por padrão. Para desligá-lo:

\begin{Verbatim}[fontsize=\small]
\usepackage[backref=false]{abnt-citacoes}
\end{Verbatim}

A opção \verb|backref| de \verb|abnt-citacoes| é repassada para
\verb|abnt-refs|, que a repassa para o biblatex.

\section{Arquivo bib}\label{sec:arquivo-bib}

\index{arquivo bib}
O arquivo \verb|.bib| segue o formato padrão do biblatex-abnt (que
é compatível com o formato BibLaTeX geral). O \verb|refs.bib| deste
guia contém exemplos de livro, artigo, recurso online, obra com
múltiplos autores e entradas duplicadas de mesmo autor/ano.

O biblatex-abnt está alinhado com a NBR 6023:2018; as diferenças
para a revisão 2025 são marginais (ISSN opcional, e-locator, etc.)
e devem ser aplicadas manualmente se exigidas pela instituição.


% =====================================================================
% CAPÍTULO 5 — TABELAS
% =====================================================================
\chapter{Tabelas}\label{ch:tabelas}

\index{tabelas}
\index{ibge-tabelas}
O pacote \verb|ibge-tabelas| implementa as Normas de Apresentação
Tabular do IBGE (3ª edição, 1993)\index{IBGE}. Ele fornece ambientes
e comandos para tabelas com três traços horizontais, sem bordas
verticais laterais, com rodapé padronizado.

\section{Ambiente tabela}\label{sec:ambiente-tabela}

O ambiente \verb|tabela{Título}| é um wrapper que integra float, caption
e rodapé como unidade coesa. Dentro dele, usa-se o ambiente interno
\verb|tabelaibge| com a opção \verb|colspec|:

\begin{Verbatim}[fontsize=\small]
\begin{tabela}{Título da tabela}
  \label{tab:exemplo}
  \ibgeFonte{Elaboração própria.}
  \begin{tabelaibge}{colspec={lrr}}
    Coluna A & Coluna B & Coluna C \\
    Dado 1   & \num{100} & \num{50.5} \\
    Dado 2   & \num{200} & \num{75.3} \\
  \end{tabelaibge}
\end{tabela}
\end{Verbatim}

\section{Rodapé da tabela}\label{sec:rodape-tabela}

\index{rodapé da tabela}
O rodapé segue ordem rígida: Fonte, Nota geral, Notas específicas.
Os comandos devem ser usados dentro do ambiente \verb|tabela|, antes
de \verb|\begin{tabelaibge}|:

\begin{alineas}
  \item \verb|\ibgeFonte{texto}| --- fonte obrigatória (seção 4.10).
    A versão com estrela \verb|\ibgeFonte*{texto}| imprime ``Fontes:''
    (plural);
  \item \verb|\ibgeNota{texto}| --- nota geral (seção 4.11).
    A versão com estrela \verb|\ibgeNota*{texto}| imprime ``Notas:''
    (plural);
  \item \verb|\ibgeNotaEsp{id}{texto}| --- nota específica (seção 4.12),
    com \verb|\ibgeChamada{id}| no corpo da tabela para inserir a
    chamada remissiva.
\end{alineas}

\section{Tabela com rodapé completo}\label{sec:tabela-completa}

A Tabela~\ref{tab:demografica} demonstra o rodapé completo com fonte,
nota geral e notas específicas.

\begin{tabela}{População residente, por região --- Brasil --- 2022}
  \label{tab:demografica}
  \ibgeFonte{IBGE, Censo Demográfico 2022.}
  \ibgeNota{Dados sujeitos a retificação.}
  \ibgeNotaEsp{1}{Inclui estimativas para municípios sem coleta}
  \ibgeNotaEsp{2}{Exclusive a população rural de Rondônia e Acre}
  \begin{tabelaibge}{colspec={lrrr}}
    Região       & Homens\ibgeChamada{1} & Mulheres & Total\ibgeChamada{2} \\
    Norte        & \num{9200000}   & \num{9100000}  & \num{18300000} \\
    Nordeste     & \num{27500000}  & \num{29000000} & \num{56500000} \\
    Sudeste      & \num{42100000}  & \num{44900000} & \num{87000000} \\
    Sul          & \num{14800000}  & \num{15200000} & \num{30000000} \\
    Centro-Oeste & \num{8200000}   & \num{8300000}  & \num{16500000} \\
  \end{tabelaibge}
\end{tabela}

\section{Sinais convencionais}\label{sec:sinais}

\index{sinais convencionais}
O pacote define comandos semânticos para os sinais convencionais do
IBGE (seção 4.8.1):

\begin{alineas}
  \item \verb|\ibgeSinalZero| --- dado igual a zero, não resultante
    de arredondamento: \ibgeSinalZero;
  \item \verb|\ibgeSinalNA| --- não se aplica: \ibgeSinalNA;
  \item \verb|\ibgeSinalND| --- dado não disponível: \ibgeSinalND;
  \item \verb|\ibgeSinalOmitido| --- dado omitido: \ibgeSinalOmitido;
  \item \verb|\ibgeSinalArrZero[casas]| --- zero por arredondamento,
    ex.: \verb|\ibgeSinalArrZero[2]| produz ``\ibgeSinalArrZero[2]'';
  \item \verb|\ibgeSinalArrZero*[casas]| --- idem, valor negativo,
    ex.: \verb|\ibgeSinalArrZero*[2]| produz ``\ibgeSinalArrZero*[2]''.
\end{alineas}

A Tabela~\ref{tab:sinais} demonstra o uso dos sinais:

\begin{tabela}{Demonstração dos sinais convencionais}
  \label{tab:sinais}
  \ibgeFonte{Dados fictícios para demonstração.}
  \begin{tabelaibge}{colspec={lrrr}}
    Indicador   & 2020              & 2021            & 2022 \\
    Produção A  & \num{1500.20}     & \num{1620.50}   & \num{1780.30} \\
    Produção B  & \ibgeSinalZero    & \num{42.10}     & \num{85.60} \\
    Produção C  & \ibgeSinalNA      & \ibgeSinalNA    & \num{10.00} \\
    Produção D  & \ibgeSinalND      & \num{320.00}    & \num{340.00} \\
    Produção E  & \ibgeSinalOmitido & \ibgeSinalOmitido & \num{5.00} \\
    Arrend.     & \ibgeSinalArrZero[2] & \num{0.10}   & \num{0.20} \\
  \end{tabelaibge}
\end{tabela}

\section{Séries temporais}\label{sec:series}

\index{séries temporais}
Comandos para apresentação de tempo (IBGE seção 5):

\begin{alineas}
  \item \verb|\ibgeSerieConsec{1981}{1985}| --- série consecutiva:
    \ibgeSerieConsec{1981}{1985};
  \item \verb|\ibgeSerieNConsec{1981}{1985}| --- série não consecutiva:
    \ibgeSerieNConsec{1981}{1985};
  \item \verb|\ibgeSafra{2020}{2021}| --- safra agrícola (recebe anos
    com 4 dígitos): \ibgeSafra{2020}{2021}.
\end{alineas}

\section{Classes de frequência}\label{sec:classes-freq}

\index{classes de frequência}
Comandos para notação textual de classes de frequência (IBGE seção 6):

\begin{alineas}
  \item \verb|\ibgeClasseIE{10}{20}| --- inclui inferior, exclui
    superior: \ibgeClasseIE{10}{20};
  \item \verb|\ibgeClasseEI{10}{20}| --- exclui inferior, inclui
    superior: \ibgeClasseEI{10}{20};
  \item \verb|\ibgeClasseII{10}{20}| --- inclui ambos:
    \ibgeClasseII{10}{20}.
\end{alineas}

\section{Nota de arredondamento}\label{sec:nota-arrend}

\index{nota de arredondamento}
O comando \verb|\ibgeNotaArrend| produz o texto padronizado:

\ibgeNotaArrend

\section{Formatação numérica}\label{sec:num}

\index{formatação numérica}
\index{siunitx}
O pacote configura o siunitx com vírgula decimal e espaço fino como
separador de milhar (padrão IBGE). O comando \verb|\num{1234567}|
produz: \num{1234567}.

\section{Opções do pacote}\label{sec:opcoes-ibge}

\index{chamada remissiva}
As opções disponíveis para \verb|ibge-tabelas| são:

\begin{alineas}
  \item \verb|chamada=parenteses| (padrão), \verb|chamada=colchetes|
    ou \verb|chamada=exponencial| --- formato da chamada remissiva;
  \item \verb|subordinar=true| --- subordina a numeração das tabelas
    aos capítulos (ex.: ``Tabela 1.1'' em vez de ``Tabela 1'').
\end{alineas}


% =====================================================================
% CAPÍTULO 6 — OUTROS ELEMENTOS
% =====================================================================
\chapter{Outros elementos}\label{ch:outros}

\section{Resumo e abstract}\label{sec:resumo-detalhe}

\index{resumo}
\index{abstract}
\index{abnt-resumo}
O pacote \verb|abnt-resumo| (NBR 6028:2021) fornece os comandos
\verb|\imprimirresumo{texto}{palavras-chave}| e
\verb|\imprimirabstract{text}{keywords}|. O fluxo de quando e onde
usá-los foi apresentado no Capítulo~\ref{ch:estrutura}. Esta seção
detalha as opções e orientações da norma.

Opções do pacote:

\begin{alineas}
  \item \verb|titulo-resumo=TEXTO| --- altera o título do resumo
    (padrão: ``RESUMO'');
  \item \verb|titulo-abstract=TEXTO| --- altera o título do abstract
    (padrão: ``ABSTRACT'').
\end{alineas}

Exemplo de uso com opções:

\begin{Verbatim}[fontsize=\small]
\usepackage[titulo-resumo=RESUMO GERAL]{abnt-resumo}
\end{Verbatim}

Orientações da NBR 6028: o resumo deve ser redigido em parágrafo único,
com extensão entre 150 e 500 palavras para trabalhos acadêmicos, e
na terceira pessoa do singular.

\section{Índice remissivo}\label{sec:indice-detalhe}

\index{índice remissivo}
\index{abnt-indice}
O pacote \verb|abnt-indice| (NBR 6034:2004) integra o imakeidx para
geração de índices remissivos com ordenação correta para o português.

\subsection{Indexação de termos}

Para indexar um termo, use o comando padrão do \LaTeX{}:

\begin{Verbatim}[fontsize=\small]
\index{termo}
\index{termo!subtermo}
\end{Verbatim}

O pacote aplica automaticamente strip de acentos para a chave de
ordenação (apenas com makeindex). Assim, \verb|\index{água}| é
convertido internamente para \verb|\index{agua@água}|.

\subsection{Remissivas}

\index{remissivas}
Dois comandos geram entradas de remissiva no índice:

\begin{alineas}
  \item \verb|\indiceVer{ONU}{Organização das Nações Unidas}| --- gera
    ``ONU, \textit{ver} Organização das Nações Unidas'';
  \item \verb|\indiceVerTambem{férias}{licença}| --- gera
    ``férias, \textit{ver também} licença''.
\end{alineas}

\indiceVer{ONU}{Organização das Nações Unidas}
\indiceVerTambem{normas}{padronização}

\subsection{Impressão do índice}

O comando \verb|\imprimirindice| gera o índice remissivo. Para índices
nomeados (múltiplos índices), use o argumento opcional:

\begin{Verbatim}[fontsize=\small]
\imprimirindice          % índice padrão
\imprimirindice[autores] % índice nomeado "autores"
\end{Verbatim}

\subsection{Opções do pacote}

\begin{alineas}
  \item \verb|titulo=TEXTO| --- altera o título do índice
    (padrão: ``ÍNDICE'');
  \item \verb|processador=makeindex| (padrão) ou
    \verb|processador=xindy| --- seleciona o processador de índice.
\end{alineas}

\section{Lombada}\label{sec:lombada-detalhe}

\index{lombada}
\index{abnt-lombada}
O pacote \verb|abnt-lombada| (NBR 12225:2023) gera uma página com o
desenho da lombada para impressão e encadernação.

\subsection{Conteúdo da lombada}

Os comandos para definir o conteúdo são:

\begin{alineas}
  \item \verb|\lombadaAutor{nome}| --- nome do autor (o pacote herda
    de \verb|\autor| automaticamente);
  \item \verb|\lombadaTitulo{título}| --- título (o pacote herda
    de \verb|\titulo| automaticamente);
  \item \verb|\lombadaElementos{v.~1}| --- elementos alfanuméricos
    (volume, fascículo);
  \item \verb|\lombadaLogomarca{arquivo.png}| --- logomarca da
    instituição (arquivo de imagem);
  \item \verb|\lombadaLogomarcaTexto{SIGLA}| --- logomarca em texto
    (alternativa ao arquivo de imagem).
\end{alineas}

O comando \verb|\imprimirlombada| gera a página com a lombada.

\subsection{Opções do pacote}

\begin{alineas}
  \item \verb|disposicao=descendente| (padrão) ou
    \verb|disposicao=horizontal| --- disposição do texto;
  \item \verb|largura=2cm| (padrão) --- largura da lombada;
  \item \verb|paginas=300| --- calcula a largura a partir do número
    de páginas (aproximação);
  \item \verb|espessura=0.1mm| (padrão) --- espessura por folha
    (usada no cálculo com \verb|paginas|).
\end{alineas}


% =====================================================================
% CAPÍTULO 7 — PERSONALIZAÇÃO
% =====================================================================
\chapter{Personalização}\label{ch:personalizacao}

\index{personalização}

\section{Links coloridos}\label{sec:links}

\index{hyperref}
\index{links}
Por padrão, o hyperref é carregado com \verb|hidelinks| (links
funcionais sem indicação visual). Para reativar a cor dos links:

\begin{Verbatim}[fontsize=\small]
\hypersetup{colorlinks, allcolors=blue}
\end{Verbatim}

Para bordas coloridas em vez de texto colorido:

\begin{Verbatim}[fontsize=\small]
\hypersetup{allbordercolors=blue}
\end{Verbatim}

\section{Maiúsculas automáticas}\label{sec:maiusculas-config}

Para desligar a conversão automática de maiúsculas nos títulos de seção
primária e secundária:

\begin{Verbatim}[fontsize=\small]
\usepackage[maiusculas=false]{abnt-numeracao}
\end{Verbatim}

\section{Sumário}\label{sec:sumario-config}

\index{sumário}
\index{abnt-sumario}
O pacote \verb|abnt-sumario| aceita duas opções:

\begin{alineas}
  \item \verb|pontilhado=false| --- remove o pontilhado entre o título
    e o número de página;
  \item \verb|titulo=CONTEÚDO| --- altera o título do sumário
    (padrão: ``SUMÁRIO'').
\end{alineas}

Exemplo:

\begin{Verbatim}[fontsize=\small]
\usepackage[pontilhado=false, titulo=CONTEÚDO]{abnt-sumario}
\end{Verbatim}

\section{Uso standalone dos pacotes}\label{sec:uso-standalone}

\index{uso standalone}
Os pacotes podem ser usados fora das classes ABNT, com
\verb|\documentclass{article}| ou \verb|\documentclass{report}|.
Nesse caso, o usuário precisa carregar manualmente os pacotes e
configurar margens, fontes e espaçamento. Exemplo:

\begin{Verbatim}[fontsize=\small]
\documentclass[a4paper, 12pt]{report}
\usepackage[T1]{fontenc}
\usepackage[brazil]{babel}
\usepackage{abnt-numeracao}
\usepackage{abnt-sumario}
\usepackage{abnt-citacoes}
\addbibresource{refs.bib}
\end{Verbatim}

As opções específicas de cada pacote estão documentadas nos respectivos
capítulos deste guia (Capítulos~\ref{ch:numeracao} a~\ref{ch:outros}).


% =====================================================================
% PÓS-TEXTUAIS
% =====================================================================
\postextual

\imprimirreferencias

% =====================================================================
% APÊNDICE A — CHECKLIST DE CONFORMIDADE ABNT
% =====================================================================
\appendix
\chapter{Checklist de conformidade ABNT}\label{ap:checklist}

\index{checklist}
A tabela a seguir lista requisitos que as normas ABNT exigem mas que
o \LaTeX{} \textbf{não} automatiza. É responsabilidade do usuário
verificar cada item.

{%
\small
\singlespacing

\noindent\textbf{Redação e estilo}

\begin{tabelaibge}{colspec={clp{3.8cm}p{5.5cm}}}
  \# & Norma & Requisito & Ação do usuário \\
  1 & NBR 6024 \S4.2 & Pontuação de alíneas &
    Letras minúsculas terminadas por ponto e vírgula, última por ponto.
    Responsabilidade do usuário. \\
  2 & NBR 6028 \S4.1.8 & Contagem de palavras do resumo &
    150--500 palavras para trabalhos acadêmicos. Não verificado
    automaticamente. \\
  3 & NBR 6028 & Tipo de resumo &
    Informativo vs indicativo. Responsabilidade do usuário. \\
  4 & NBR 6028 \S4.1.4 & Verbo na terceira pessoa &
    Estilo de redação do resumo. Responsabilidade do usuário. \\
\end{tabelaibge}

\vspace{12pt}
\noindent\textbf{Elementos do documento}

\begin{tabelaibge}{colspec={clp{3.8cm}p{5.5cm}}}
  \# & Norma & Requisito & Ação do usuário \\
  5 & NBR 14724 & Ordem dos elementos pré-textuais &
    A classe fornece os comandos mas não impõe sequência. Seguir a
    ordem do Capítulo~\ref{ch:estrutura}. \\
  6 & NBR 14724 & Folha de aprovação &
    Inserida após a defesa. Antes da defesa, usar placeholders. \\
  7 & NBR 14724 & Errata &
    Elemento opcional não implementado; inserir manualmente se
    necessário. \\
  8 & NBR 14724 & Listas de ilustrações, tabelas, siglas, símbolos &
    Não automatizadas pelas classes. Usar
    \code{\bsl listoffigures},
    \code{\bsl listoftables} ou construir manualmente. \\
  9 & NBR 14724 & Ficha catalográfica &
    Verso da folha de rosto. Geralmente fornecida pela biblioteca;
    inserir como PDF. \\
\end{tabelaibge}

\vspace{12pt}
\noindent\textbf{Tabelas (IBGE)}

\begin{tabelaibge}{colspec={clp{3.8cm}p{5.5cm}}}
  \# & Norma & Requisito & Ação do usuário \\
  10 & IBGE \S4.10 & Fonte obrigatória em tabelas &
    Warning emitido, mas o usuário deve incluir \code{\bsl ibgeFonte}. \\
  11 & IBGE \S7.4b & Correção por razão de arredondamento &
    Requer dados brutos; não automatizável. \\
  12 & IBGE \S8.1--8.2 & Diagramação de tabelas largas &
    Responsabilidade do usuário. \\
  13 & IBGE & Tabelas longas &
    \code{longtblr} do tabularray está disponível mas não integrado ao
    ambiente \code{tabela}. \\
  14 & IBGE \S6.1b--6.3b & Notação gráfica de classes de frequência &
    Apenas notação textual implementada. \\
  15 & IBGE/siunitx & Coluna S do siunitx com tabularray &
    Incompatível em versões anteriores; usar \code{\bsl num\{\}}. \\
\end{tabelaibge}

\vspace{12pt}
\noindent\textbf{Índice}

\begin{tabelaibge}{colspec={clp{3.8cm}p{5.5cm}}}
  \# & Norma & Requisito & Ação do usuário \\
  16 & NBR 6034 \S6.8 & Marca ``(continuação)'' &
    Em índices que quebram de página. Não implementado. \\
  17 & NBR 6034 \S6.5 & Letras iniciais no running header &
    Não implementado. \\
\end{tabelaibge}

\vspace{12pt}
\noindent\textbf{Sumário}

\begin{tabelaibge}{colspec={clp{3.8cm}p{5.5cm}}}
  \# & Norma & Requisito & Ação do usuário \\
  18 & NBR 6027 \S5.3 & Obras coletivas &
    Sumário com vários autores. Não implementado (caso raro). \\
  19 & NBR 6027 \S5.5 & Traduções &
    Sumário traduzido. Não implementado (caso raro). \\
  20 & NBR 6027 \S5.6 & Sumário por idioma &
    Não implementado (caso raro). \\
  21 & NBR 6027 & Largura fixa do indicativo &
    O indicativo usa largura fixa de 4em. Ajustar manualmente com
    código expl3 se necessário. \\
\end{tabelaibge}

\vspace{12pt}
\noindent\textbf{Lombada}

\begin{tabelaibge}{colspec={clp{3.8cm}p{5.5cm}}}
  \# & Norma & Requisito & Ação do usuário \\
  22 & NBR 12225 \S4.1c & Modo misto &
    Combinação descendente/horizontal. Não implementado. \\
  23 & NBR 12225 \S4.3.3 & Título de margem de capa &
    Fora do escopo do pacote. \\
\end{tabelaibge}

\vspace{12pt}
\noindent\textbf{Citações}

\begin{tabelaibge}{colspec={clp{3.8cm}p{5.5cm}}}
  \# & Norma & Requisito & Ação do usuário \\
  24 & NBR 10520 \S7.1.1 & Detecção automática de citação longa &
    Citação com mais de 3 linhas deve usar \code{citacaodireta}.
    A detecção não é viável em TeX; o usuário decide. \\
\end{tabelaibge}

\vspace{12pt}
\noindent\textbf{Referências}

\begin{tabelaibge}{colspec={clp{3.8cm}p{5.5cm}}}
  \# & Norma & Requisito & Ação do usuário \\
  25 & NBR 6023 & Diferenças NBR 6023:2018 vs 2025 &
    biblatex-abnt alinhado com 2018. Diferenças (ISSN opcional,
    e-locator, etc.) devem ser aplicadas manualmente. \\
\end{tabelaibge}

}% end of \small\singlespacing group


% =====================================================================
% APÊNDICE B — REFERÊNCIA RÁPIDA DE COMANDOS
% =====================================================================
\chapter{Referência rápida de comandos}\label{ap:referencia}

\index{referência rápida}
A tabela a seguir lista todos os comandos públicos do projeto, organizados
por pacote ou classe.

{%
\small
\singlespacing

\noindent\textbf{abnt.cls --- Classe base}

\begin{tabelaibge}{colspec={lp{3.5cm}p{5.5cm}}}
  Comando & Argumentos & Descrição \\
  \code{\bsl titulo}          & \code{\{texto\}}   & Define o título do trabalho \\
  \code{\bsl subtitulo}       & \code{\{texto\}}   & Define o subtítulo \\
  \code{\bsl autor}           & \code{\{nome\}}    & Define o nome do autor \\
  \code{\bsl orientador}      & \code{\{nome\}}    & Define o nome do orientador \\
  \code{\bsl coorientador}    & \code{\{nome\}}    & Define o nome do coorientador \\
  \code{\bsl instituicao}     & \code{\{nome\}}    & Define a instituição \\
  \code{\bsl local}           & \code{\{cidade\}}  & Define a cidade \\
  \code{\bsl ano}             & \code{\{ano\}}     & Define o ano \\
  \code{\bsl natureza}        & \code{\{texto\}}   & Define a natureza do trabalho \\
  \code{\bsl programa}        & \code{\{texto\}}   & Define o programa \\
  \code{\bsl area}            & \code{\{texto\}}   & Define a área de concentração \\
  \code{\bsl imprimircapa}    & ---              & Gera a capa \\
  \code{\bsl imprimirfolhaderosto} & ---         & Gera a folha de rosto \\
  \code{\bsl pretextual}      & ---              & Inicia modo pré-textual \\
  \code{\bsl textual}         & ---              & Inicia modo textual \\
\end{tabelaibge}

\vspace{12pt}
\noindent\textbf{abnt-trabalho-academico.cls --- Trabalho acadêmico}

\begin{tabelaibge}{colspec={lp{3.5cm}p{5.5cm}}}
  Comando & Argumentos & Descrição \\
  \code{\bsl dataaprovacao}   & \code{\{data\}} & Define data de aprovação \\
  \code{\bsl membrodabanca}   & \code{\{nome\}\{tit.\}} & Adiciona membro da banca \\
  \code{\bsl imprimirfolhadeaprovacao} & --- & Gera folha de aprovação \\
  \code{\bsl imprimirdedicatoria} & \code{\{texto\}} & Gera dedicatória \\
  \code{\bsl imprimiragradecimentos} & \code{\{texto\}} & Gera agradecimentos \\
  \code{\bsl imprimirepigrafe} & \code{\{texto\}} & Gera epígrafe \\
  \code{\bsl postextual}      & ---              & Inicia modo pós-textual \\
\end{tabelaibge}

\vspace{12pt}
\noindent\textbf{abnt-numeracao.sty --- Numeração progressiva (NBR 6024)}

\begin{tabelaibge}{colspec={lp{3.5cm}p{5.5cm}}}
  Comando/Ambiente & Argumentos/Opções & Descrição \\
  \code{alineas}       & (ambiente)        & Lista com alíneas (a, b, c\ldots) \\
  \code{subalineas}    & (ambiente)        & Sublista com travessão \\
  opção \code{maiusculas} & \code{true/false} & Ativa/desativa maiúsculas \\
\end{tabelaibge}

\vspace{12pt}
\noindent\textbf{abnt-sumario.sty --- Sumário (NBR 6027)}

\begin{tabelaibge}{colspec={lp{3.5cm}p{5.5cm}}}
  Comando & Argumentos/Opções & Descrição \\
  \code{\bsl imprimirsumario} & --- & Gera o sumário ABNT \\
  opção \code{titulo}     & \code{TEXTO} & Altera título do sumário \\
  opção \code{pontilhado} & \code{true/false} & Ativa/desativa pontilhado \\
\end{tabelaibge}

\vspace{12pt}
\noindent\textbf{abnt-resumo.sty --- Resumo e abstract (NBR 6028)}

\begin{tabelaibge}{colspec={lp{3.5cm}p{5.5cm}}}
  Comando & Argumentos & Descrição \\
  \code{\bsl imprimirresumo}   & \code{\{texto\}\{palavras\}} & Gera o resumo \\
  \code{\bsl imprimirabstract}  & \code{\{text\}\{keywords\}}  & Gera o abstract \\
  opção \code{titulo-resumo} & \code{TEXTO} & Altera título do resumo \\
  opção \code{titulo-abstract} & \code{TEXTO} & Altera título do abstract \\
\end{tabelaibge}

\vspace{12pt}
\noindent\textbf{abnt-citacoes.sty --- Citações (NBR 10520)}

\begin{tabelaibge}{colspec={lp{3.5cm}p{5.5cm}}}
  Comando/Ambiente & Argumentos/Opções & Descrição \\
  \code{citacaodireta}  & (ambiente)    & Citação direta longa \\
  opção \code{sistema}  & \code{autor-data/numerico} & Sistema de citação \\
  opção \code{backref}  & \code{true/false} & Ativa/desativa backref \\
\end{tabelaibge}

\vspace{12pt}
\noindent\textbf{abnt-refs.sty --- Referências (NBR 6023)}

\begin{tabelaibge}{colspec={lp{3.5cm}p{5.5cm}}}
  Comando & Argumentos & Descrição \\
  \code{\bsl imprimirreferencias} & --- & Gera lista de referências \\
  opção \code{backref} & \code{true/false} & Ativa/desativa backref \\
\end{tabelaibge}

\vspace{12pt}
\noindent\textbf{ibge-tabelas.sty --- Tabelas (IBGE 1993)}

\begin{tabelaibge}{colspec={lp{3.5cm}p{5.5cm}}}
  Comando/Ambiente & Argumentos & Descrição \\
  \code{tabela}           & \code{\{título\}} (ambiente) & Wrapper: float + caption + rodapé \\
  \code{tabelaibge}       & \code{\{colspec=\{...\}\}} (amb.) & Tabela interna com estilo IBGE \\
  \code{\bsl ibgeFonte}       & \code{\{texto\}}    & Fonte no rodapé (obrigatória) \\
  \code{\bsl ibgeFonte*}      & \code{\{texto\}}    & Fontes (plural) \\
  \code{\bsl ibgeNota}        & \code{\{texto\}}    & Nota geral \\
  \code{\bsl ibgeNota*}       & \code{\{texto\}}    & Notas (plural) \\
  \code{\bsl ibgeNotaEsp}     & \code{\{id\}\{texto\}} & Nota específica \\
  \code{\bsl ibgeChamada}     & \code{\{id\}}       & Chamada remissiva no corpo \\
  \code{\bsl ibgeSinalZero}   & ---               & Sinal: dado = zero \\
  \code{\bsl ibgeSinalNA}     & ---               & Sinal: não se aplica \\
  \code{\bsl ibgeSinalND}     & ---               & Sinal: não disponível \\
  \code{\bsl ibgeSinalOmitido} & ---              & Sinal: dado omitido \\
  \code{\bsl ibgeSinalArrZero} & \code{[casas]}   & Zero por arredondamento \\
  \code{\bsl ibgeSinalArrZero*} & \code{[casas]}  & Idem, valor negativo \\
  \code{\bsl ibgeSerieConsec} & \code{\{a1\}\{a2\}}   & Série consecutiva \\
  \code{\bsl ibgeSerieNConsec} & \code{\{a1\}\{a2\}}  & Série não consecutiva \\
  \code{\bsl ibgeSafra}       & \code{\{a1\}\{a2\}}   & Safra agrícola \\
  \code{\bsl ibgeClasseIE}    & \code{\{min\}\{max\}} & Classe: inclui inf., exclui sup. \\
  \code{\bsl ibgeClasseEI}    & \code{\{min\}\{max\}} & Classe: exclui inf., inclui sup. \\
  \code{\bsl ibgeClasseII}    & \code{\{min\}\{max\}} & Classe: inclui ambos \\
  \code{\bsl ibgeNotaArrend}  & ---               & Nota de arredondamento \\
  \code{\bsl num}             & \code{\{número\}}   & Formatação numérica (siunitx) \\
  opção \code{chamada}    & \code{par./colch./exp.} & Formato da chamada \\
  opção \code{subordinar} & \code{true/false} & Subordina numeração a capítulos \\
\end{tabelaibge}

\vspace{12pt}
\noindent\textbf{abnt-indice.sty --- Índice remissivo (NBR 6034)}

\begin{tabelaibge}{colspec={lp{3.5cm}p{5.5cm}}}
  Comando & Argumentos & Descrição \\
  \code{\bsl index}           & \code{\{termo\}}       & Indexa um termo (padrão LaTeX) \\
  \code{\bsl indiceVer}       & \code{\{de\}\{para\}}    & Remissiva ``ver'' \\
  \code{\bsl indiceVerTambem}  & \code{\{de\}\{para\}}    & Remissiva ``ver também'' \\
  \code{\bsl imprimirindice}  & \code{[nome]}        & Gera o índice remissivo \\
  opção \code{titulo}     & \code{TEXTO}         & Altera título do índice \\
  opção \code{processador} & \code{makeindex/xindy} & Seleciona processador \\
\end{tabelaibge}

\vspace{12pt}
\noindent\textbf{abnt-lombada.sty --- Lombada (NBR 12225)}

\begin{tabelaibge}{colspec={lp{3.5cm}p{5.5cm}}}
  Comando & Argumentos & Descrição \\
  \code{\bsl lombadaAutor}     & \code{\{nome\}}    & Define autor na lombada \\
  \code{\bsl lombadaTitulo}    & \code{\{título\}}  & Define título na lombada \\
  \code{\bsl lombadaElementos} & \code{\{texto\}}   & Elementos alfanuméricos \\
  \code{\bsl lombadaLogomarca} & \code{\{arquivo\}} & Logomarca (imagem) \\
  \code{\bsl lombadaLogomarcaTexto} & \code{\{texto\}} & Logomarca (texto) \\
  \code{\bsl imprimirlombada}  & ---              & Gera a página da lombada \\
  opção \code{disposicao}  & \code{desc./horiz.} & Disposição do texto \\
  opção \code{largura}     & \code{dimensão}  & Largura da lombada \\
  opção \code{paginas}     & \code{número}    & Calcula largura por páginas \\
  opção \code{espessura}   & \code{dimensão}  & Espessura por folha \\
\end{tabelaibge}

}% end of \small\singlespacing group


% =====================================================================
% ÍNDICE REMISSIVO E LOMBADA
% =====================================================================
\imprimirindice

\imprimirlombada

\end{document}
